%%%%%%%%%%%%%%%%%%%%%%%%%%%%%%%%%%%%%%%%%%%%%%%%%%%%%%%
% % Lines starting with % are comments, which are ignored.
% % This is a handy way of indicating the date and version of
% % your document, to wit:
% %
% % LaTeX sample file
% % Modified March, 2002
% %
%%%%%%%%%%%%%%%%%%%%%%%%%%%%%%%%%%%%%%%%%%%%%%%%%%%%%%%
\DocumentMetadata{
  % PDF standards: ONE key whose value is a comma-separated list --
  % UA-2 (accessibility) plus A-4f (archival; the archival companion of
  % UA-2). Do NOT use separate pdfstandard=... keys: with ordinary
  % key-value parsing only the last one takes effect.
  pdfstandard={UA-2,A-4f},
  lang=en-US,
  tagging=on,
  tagging-setup={
    math/setup=mathml-SE
  }
}

\documentclass{ltx-talk}
% lecture-macros.tex -- shared package setup for lecture-template.tex
% (ltx-talk slides / handouts).
%
% NOTE ON PATHS: every template in this repository is stored one level up
% from where it is compiled (see README.md). All relative \input paths
% below are resolved against the folder the document is compiled in,
% which sits one directory above this file.

% --- Tables, bookmarks, strikethrough ------------------------------------
% booktabs: professional table rules (\toprule, \midrule, \bottomrule)
% bookmark: PDF outline entries (PDF bookmarks)
% ulem[normalem]: \sout strikethrough that survives math mode
\usepackage{booktabs}
\usepackage{bookmark}
\usepackage[normalem]{ulem}
\usepackage{multicol}

% --- Math ------------------------------------------------------------
% amsmath: equation environments and numbering
% mathrsfs: \mathfrak  |  stmaryrd: extra symbols
\usepackage{amsmath}
\usepackage{mathrsfs}
\usepackage{stmaryrd}

% --- Layout -------------------------------------------------------------
% float: [H] float placement
% fontspec: Unicode/OpenType fonts (compile with XeLaTeX or LuaLaTeX)
\usepackage{float}
\usepackage{fontspec}

% --- Colors ---------------------------------------------------------------
% dvipsnames: extended named color palette (the Accessible* colors below
% are built on it);
% table: \rowcolors / \cellcolor for shading tables.
% NOTE: xcolor may only be loaded ONCE per document -- a second
% \usepackage{xcolor} with different options is a fatal "Option clash"
% error, so ALL of its options are listed here.
\usepackage[dvipsnames,table]{xcolor}

% --- Multimedia -----------------------------------------------------------
% animate: embeds audio/video (and legacy Flash) media in a frame.
\usepackage{animate}

% --- Bibliography ----------------------------------------------------------
% biblatex (biber backend): citations for lecture notes.
\usepackage[
  backend=biber,
  style=alphabetic,
  url=true,
  doi=true,
  isbn=true,
  backref=true
]{biblatex}

% --- Colored boxes ----------------------------------------------------------
% talkthemetcolorbox: the colored-box theme for ltx-talk slides
% (shadow/rounded/shaded blocks, showtitle, titlepage styles).
\usepackage[
  shadow,
  rounded,
  shaded,
  blocks,
  showtitle,
  titlepage
]{talkthemetcolorbox}

% --- "theindex" helper environment ------------------------------------------
% Provides a lightweight theindex environment (an itemize) and the
% \@idxitem/\subitem/\subsubitem/\indexspace macros that LaTeX's standard
% theindex machinery expects, so hand-written index pages indent correctly.
\makeatletter
\newenvironment{theindex}
  {\begin{itemize}}
  {\end{itemize}}
\providecommand\@idxitem{\par\hangindent 40\p@}
\providecommand\subitem{\@idxitem \hspace*{20\p@}}
\providecommand\subsubitem{\@idxitem \hspace*{30\p@}}
\providecommand\indexspace{\par \vskip 10\p@ \@plus5\p@ \@minus3\p@\relax}
\makeatother

% --- Index -----------------------------------------------------------------
% imakeidx: builds the printed index.
% The redefinition below appends "|hyperpage" to every index entry so that,
% in the PDF, each index entry is a link to its page.
\usepackage{imakeidx}
\makeatletter
\let\imki@orig@wrindexentrysplit\imki@wrindexentrysplit
\renewcommand\imki@wrindexentrysplit[3]{%
  \imki@orig@wrindexentrysplit{#1}{#2|hyperpage}{#3}%
}

\makeatletter
% --- Section titles: capture + one bookmark per logical section -------------
% \section is redefined so that:
%   1. the section title is captured into the plain-TeX register
%      \mytalksectitletoks (reprinted verbatim on the auto divider frame,
%      see the cmd/section/after hook below -- plain-TeX storage is
%      catcode-safe, unlike expanding the title in a hook), and
%   2. exactly ONE depth-1 PDF bookmark is created per \section call
%      (hyperref's own bookmark is suppressed via \Hy@writebookmark).
\DeclareCommandCopy{\OriginalSectionForBookmarks}{\section}
\newtoks\mytalksectitletoks

\RenewDocumentCommand{\section}{s o m}{%
  \begingroup
    \def\Hy@writebookmark##1##2##3##4##5{}%
    \global\mytalksectitletoks={#3}%
    \IfBooleanTF{#1}
      {\phantomsection\IfNoValueTF{#2}{\OriginalSectionForBookmarks*{#3}}{\OriginalSectionForBookmarks*[#2]{#3}}}%
      {\IfNoValueTF{#2}{\OriginalSectionForBookmarks{#3}}{\OriginalSectionForBookmarks[#2]{#3}}}%
  \endgroup
  \IfBooleanTF{#1}
    {\IfNoValueTF{#2}{\pdfbookmark[1]{#3}{section-star.\thepage.\number\inputlineno}}{\pdfbookmark[1]{#2}{section-star.\thepage.\number\inputlineno}}}%
    {\IfNoValueTF{#2}{\pdfbookmark[1]{Section \thesection: #3}{section.\thesection}}{\pdfbookmark[1]{Section \thesection: #2}{section.\thesection}}}%
}
\makeatother

% --- One PDF bookmark per logical frame (NOT per overlay) ---------------------
% ltx-talk writes one PDF page per overlay, so a naive frametitle hook would
% register a new bookmark for every overlay of the same frame. The
% \g_slideutils_frame_bookmark_int register remembers the last frame that
% got a bookmark, so subsequent overlays of the same frame are skipped.
% This is also what keeps BeamerPresenter's cumulative annotations working:
% one logical frame must stay one logical unit in the PDF outline.
\ExplSyntaxOn
\int_new:N \g_slideutils_frame_bookmark_int

\AddToHook{cmd/frametitle/after}{%
  \tl_if_empty:NF \g__talk_frame_title_tl
    {
      \ifnum\value{frame}=\g_slideutils_frame_bookmark_int\relax
      \else
        \int_gset:Nn \g_slideutils_frame_bookmark_int { \value{frame} }%
        \pdfbookmark[2]{\tl_use:N \g__talk_frame_title_tl}
          {frame.\int_use:N \g_slideutils_frame_bookmark_int}
      \fi
    }
}
\ExplSyntaxOff
\makeatother

% --- Printed index + theme for the colored boxes -------------------------------
% \makeindex[...]: two-column index titled "Index", listed in the TOC.
% The two EditInstance calls style the standard header/footer instances
% of talkthemetcolorbox (cyan background; footer order title/author/
% date/framenumber).
\makeindex[columns=2, title=Index, intoc]\EditInstance{header}{std}{
  background-color = cyan,
  color            = black,
  height           = 0.75cm
}
\EditInstance{footer}{std}{
  background-color = cyan,
  element-order = {title, author, date, framenumber},
}

% --- Auto section divider frame ---------------------------------------------------
% After each \section, insert a centered, large, structure-colored frame
% that prints the section title (from \mytalksectitletoks).
\AddToHook{cmd/section/after}{%
  \begin{frame}
    \centering
    \color{structure}
    \LARGE
    \the\mytalksectitletoks
  \end{frame}%
}

% --- Tagging sockets (tagged PDF structure) ---------------------------------------
% ltx-talk exposes tagging "sockets" where the PDF structure can attach
% section titles; register one and leave it unassigned (none) so the
% class's own tagging keeps control.
\ExplSyntaxOn
\AtBeginDocument{
  \NewTaggingSocketPlug { talk / sec / title } { none } { }
  \AssignTaggingSocketPlug { talk / sec / title } { none }
}
\ExplSyntaxOff

% --- Accessibility roles ----------------------------------------------------------------
% Map \frametitle to the H1 heading role in the tagged PDF, so screen
% readers treat slide titles as headings.
\tagpdfsetup{role / new-tag = frametitle / H1}

\makeatletter
% --- Page label = frame number -------------------------------------------------------
% Redefine what feeds the PDF page label to track the frame counter,
% not the physical page counter. Replace \theframenumber below with
% whatever internal macro/property ltx-talk actually uses to produce
% the "framenumber" footer element — confirm via ltx-talk-frame.dtx.
\renewcommand\thepage{\arabic{frame}}
\makeatother

\addbibresource{refs.bib}

% --- "Problem Time!" slides ------------------------------------------------------------------
% \ProblemSlide{problem}{solution} emits two frames: the problem, then
% its solution after a \pause, each numbered by the problem counter.
\newcounter{problemcounter}
\setcounter{problemcounter}{0}
\newcommand{\ProblemSlide}[2]{
  \stepcounter{problemcounter}
  \begin{frame}
    \frametitle{Problem Time!}
    Problem \theproblemcounter: #1
  \end{frame}

  \begin{frame}
    \frametitle{Problem Time!}
    Solution to Problem \theproblemcounter: \pause

    #2
  \end{frame}
}
% \highlight[<colour>]{<stuff>}
% from https://tex.stackexchange.com/questions/33401/a-version-of-colorbox-that-works-inside-math-environments
\newcommand{\highlight}[2][yellow]{\mathchoice%
  {\colorbox{#1}{$\displaystyle#2$}}%
  {\colorbox{#1}{$\textstyle#2$}}%
  {\colorbox{#1}{$\scriptstyle#2$}}%
  {\colorbox{#1}{$\scriptscriptstyle#2$}}}%

% --- Fonts ------------------------------------------------------------------------------
% Lato (text) + LeteSansMath (math): open, accessible typefaces.
\setmainfont{Lato}[
  Ligatures=TeX,
  Numbers={Lining,Monospaced}
]
\setsansfont{Lato}[
  Ligatures=TeX,
  Numbers={Lining,Monospaced}
]
\setmathfont{LeteSansMath.otf}\usepackage{emoji}

% --- Accessible colors ---------------------------------------------------------------------
% WCAG-contrast-safe variants of red/green for text on white:
% red!80!black and green!40!black both clear the 4.5:1 requirement.
% AltYellow: soft background for highlight boxes.
\colorlet{AccessibleRed}{red!80!black}
\colorlet{AccessibleGreen}{green!40!black}
\definecolor{AltYellow}{RGB}{255,245,170}

% --- Links ---------------------------------------------------------------------------------
% hyperref + link colors: every link uses AccessibleGreen (7.2:1 on
% white), so internal, citation, and URL links all pass WCAG 1.4.3.
\usepackage{hyperref}
\hypersetup{
  colorlinks=true,
  linkcolor=AccessibleGreen,
  urlcolor=AccessibleGreen,
  citecolor=AccessibleGreen,
}

% --- Shared course macros ---------------------------------------------------------------------
% Root macros.tex: math shorthands used across all course documents
% (\RR, \NN, \ZZ, \QQ, \grad, \curl, \tr, \der, \pder, ...).
\newcommand{\RR}{\mathbb{R}}
\newcommand{\NN}{\mathbb{N}}
\newcommand{\ZZ}{\mathbb{Z}}
\newcommand{\QQ}{\mathbb{Q}}
\newcommand{\vv}{\vec{v}}
\newcommand{\ww}{\vec{w}}
\newcommand{\xx}{\vec{x}}
\newcommand{\yy}{\vec{y}}
\newcommand{\zz}{\vec{z}}
\newcommand{\uu}{\vec{u}}
\newcommand{\basis}[1]{\vec{e_{#1}}}
\newcommand{\id}{\mathrm{id}}
\newcommand{\CC}{\mathbb{C}}
\renewcommand{\div}{\operatorname{div}}
\newcommand{\grad}{\operatorname{grad}}
\newcommand{\curl}{\operatorname{curl}}
\newcommand{\tr}{\operatorname{tr}}
\newcommand{\der}[2]{\dfrac{d #1}{d #2}}
\newcommand{\pder}[3][]{\dfrac{\partial^{#1} #2}{\partial #3}}


% % Title and author(s)
%%%%%%%%%%%%%%%%%%%%%%%%%%%%%%%%%%%%%%%%%%%%%%%%%%%%%%%
\title{Lecture 3 - Limits II}
\author{Garrett Credi}
\date{09/14/2026}
%%%%%%%%%%%%%%%%%%%%%%%%%%%%%%%%%%%%%%%%%%%%%%%%%%%%%%%
% %
% % The next command allows your in import encapsulated
% % postscript files, .epsf or .eps files, which
% % contain vector graphic image data.
% %
%%%%%%%%%%%%%%%%%%%%%%%%%%%%%%%%%%%%%%%%%%%%%%%%%%%%%%%
% %   The first thanks indicates your affiliation
% %
% %  Just the name here.
% %
% % Your mailing address goes at the end.
% %
% % \thanks is also how you indicate grant support
% %
%%%%%%%%%%%%%%%%%%%%%%%%%%%%%%%%%%%%%%%%%%%%%%%%%%%%%%%

%%%%%%%%%%%%%%%%%%%%%%%%%%%%%%%%%%%%%%%%%%%%%%%%%%%%%%%
% $ The next command ensures the margins are 1 inch on
% % all sides.
% %
%%%%%%%%%%%%%%%%%%%%%%%%%%%%%%%%%%%%%%%%%%%%%%%%%%%%%%%
\begin{document}
\begin{frame}
  \maketitle
\end{frame}

\begin{frame}
  \frametitle{Announcements}

  Homework Policy: \pause

  \begin{itemize}
    \item For homework assigned on Tueday, the due date will be Thursday @ 6PM  \pause
    \item For homework assigned on Thursday, the due date will be Saturday @ 6PM
  \end{itemize} \pause

  Office Hours: \pause
  \begin{itemize}
    \item Wednesday 12:15 - 1:15 PM
    \item Thursday 4:30 - 5:30 PM \pause
  \end{itemize}

  Both are in Vincent Hall 105.
\end{frame}

\begin{frame}
  \frametitle{Last Time:}

  Last lecture we: \pause
  \begin{enumerate}
    \item Defined the notion of a \textit{limit} of a function: \pause $\lim_{x \rightarrow a} f(x) = L \in \RR$
      if we can make $f(x) \approx L$ up to some error ($\epsilon > 0$) just by forcing $x \approx a$ up to some error ($\delta > 0$) \pause
    \item Said that if no such value $L$ existed we said that $\lim_{x \rightarrow a} f(x)$ \textit{did not exist} D.N.E. \pause
    \item Computed limits graphically: \pause 
      \begin{itemize}
        \item points where the graph of a function ``broke" were points where the limit did not exist  \pause
        \item points where the graph of a function had a ``hole" had the limit exist but it \textit{\textbf{was not equal to the function's value}} \pause
        \item points where the graph was just an unbroken curve had the limit exist and \textit{equal} to the value at that point \pause
      \end{itemize}
    \item \textit{Approximated} limits using a table of values \pause
    \item Computed limits by algebraic manipulation
  \end{enumerate}
\end{frame}

\begin{frame}
  \frametitle{Today's Goals:}

  Today we will develop our ability to algebraically compute limits by: \pause
  \begin{enumerate}
    \item Studying the ``simplest" functions to take limits of, \textit{continuous functions}
    \item Discovering \textit{limit properties} that help us reduce complicated limits to easier limits \pause
    \item Combining the above to get the process of algebraically computing limits
  \end{enumerate}

  And we will also study cases where the limit is infinite.
\end{frame}

\section{Continuity}

\begin{frame}
  \frametitle{Continuous Functions I}

  To start off: let's consider the function: 
  \begin{equation}
    f(x) = \dfrac{x^3 - x^2 - x + 1}{x^2 + x - 2}
  \end{equation} \pause

  \begin{enumerate}
    \item What is the domain of $f$? \pause
    \item Using the graph of $f(x)$, what are:
      \begin{enumerate}
        \item $\lim_{x \rightarrow 3}f(x)$?
        \item $\lim_{x \rightarrow 5}f(x)$?
        \item $\lim_{x \rightarrow 1} f(x)$? \pause
      \end{enumerate}

  \end{enumerate}
\end{frame}

\begin{frame}
  \frametitle{Continuous Functions II}

  Those first two points, where we have
  \begin{equation*}
    f(3) = \lim_{x \rightarrow 3} f(x) \quad f(5) = \lim_{x \rightarrow 5} f(x)
  \end{equation*}

  are a special kind of point: \pause

  \begin{block}{Continuity at a point}\index{continuity!at a point}
    Given a function $f(x)$ and a real number $a$ , \pause
    we say that \textit{$f(x)$ is continuous at $a$} if: \pause
    \begin{enumerate}
      \item The limit $\lim_{x \rightarrow a} f(x)$ exists \pause
      \item $a$ is in the domain of $f$ \pause
      \item $\lim_{x \rightarrow a} f(x) = f(a)$
    \end{enumerate}
  \end{block}
\end{frame}

\begin{frame}
  \frametitle{Continuous Functions III}

  We can define classes of functions that are specified to be continuous at important collections of points: \pause

  \begin{block}{Continuous on $[a,b]$}\index{continuity!on an interval}
    A function $f(x)$ is said to be continuous on the interval $[a,b]$
    if, \pause for every real number $a \leq x_0 \leq b$ \pause
    we have that $f(x)$ is continuous at $x_0$. \pause
  \end{block}

  \begin{block}{Continuous Function}\index{continuity!continuous function}\index{continuous function}
    A function $f(x)$ is said to be continuous if, \pause for every real number $x_0$ in the domain of $f$, \pause
    we have that $f(x)$ is continuous at $x_0$.
  \end{block}
\end{frame}

\begin{frame}
  \frametitle{Continuous Functions IV}

  Continuous functions are the ``atoms" we want to consider when computing limits \pause
  since we can replace the difficult task of \textit{proving} what number a functions value approaches \pause
  by the easier task of \textit{evaluating} the function. \pause

  What are some continuous function? \pause

  \begin{alertblock}{Important Continuous Functions}\index{continuous function!examples of}
    The following are examples of continuous functions: \pause
    \begin{enumerate}
      \item The constant function $f_c(x) = c$ for some constant number $c$. \pause
      \item The identity function $f_{id}(x) = x$ \pause
      \item The exponential and logarithm functions $\mathrm{exp}(x) = e^x, \ln(x)$ \pause
      \item All trigonometric functions, $\sin(x), \cos(x), \tan(x)...$ \pause
    \end{enumerate}
  \end{alertblock}
\end{frame}

\begin{frame}
  \frametitle{Continuous Functions V}
  What about other familiar functions like: \pause
  \begin{enumerate}
    \item Power functions $x^n$, \pause
    \item Polynomials $c_0x^n + c_1x^{n-1} + \ldots + c_n x + c_{n+1}$, \pause
    \item Rational functions $\dfrac{p(x)}{q(x)}$ \pause
    \item Weird functions like $e^{\sin(x)} + \cos(x^n)\sqrt{x}$?
  \end{enumerate}

  We can determine the points at which these functions are continuous \pause (and thus the points at which it is easy to calculate limits) \pause
  by discovering \textit{limit properties}!
\end{frame}

\section{Limit Properties}

\begin{frame}
  \frametitle{Limit Properties: Additivity I}

  Say we have two functions $f(x)$ and $g(x)$, and we know that:
  \begin{enumerate}
    \item As $x \rightarrow 0$, $f(x) \rightarrow 1$ and \pause
    \item As $x \rightarrow 0$, $g(x) \rightarrow 0$ \pause
  \end{enumerate}

  When $x \rightarrow 0$, towards which value does the function $h(x) = f(x) + g(x)$ approach?

  \begin{enumerate}
    \item Property 1 says that as long as $x \approx 0$ then $f(x) \approx 1$ \pause
    \item Property 2 says that as long as $x \approx 0$ then $g(x) \approx 0$ \pause
  \end{enumerate}

  So if $x \approx 0$ then $h(x) = f(x) + g(x) \approx 1 + 0 = 1$!
\end{frame}

\begin{frame}
  \frametitle{Limit Properties: Additivity II}

  Let's codify this into a property: \pause

  \begin{alertblock}{Additivity of Limits}\index{limits!additivity}
    Suppose that $f(x)$ and $g(x)$ are two functions and $a$ is a real number so that \pause
    \begin{enumerate}
      \item $\lim_{x \rightarrow a} f(x)$ exists \pause and is equal to $b$ \pause
      \item $\lim_{x \rightarrow a} g(x)$ exists \pause and is equal to $c$ \pause
    \end{enumerate}

    Then we have that $\lim_{x \rightarrow a} (f(x) + g(x))$ exists \pause and is equal to $b + c$.
  \end{alertblock}
\end{frame}

\begin{frame}
  \frametitle{Limit Properties: Multiplicativity}

  \begin{alertblock}{Multiplicativity of Limits}\index{limits!multiplicativity}
    Suppose that $f(x)$ and $g(x)$ are two functions and $a$ is a real number so that \pause
    \begin{enumerate}
      \item $\lim_{x \rightarrow a} f(x)$ exists \pause and is equal to $b$ \pause
      \item $\lim_{x \rightarrow a} g(x)$ exists \pause and is equal to $c$ \pause
    \end{enumerate}

    Then we have that $\lim_{x \rightarrow a} (f(x)g(x))$ exists \pause and is equal to $bc$.
  \end{alertblock}

  Similarly: \pause

  \begin{alertblock}{Limits of Multiplicative Inverses} 
    Suppose that $f(x)$ is a function and $a$ is a real number so that \pause
    $\lim_{x \rightarrow a} f(x)$ exists, \pause is equal to some real number $b$, \pause
    and that $b \neq 0$.

    Then we have that $\lim_{x \rightarrow a} \dfrac{1}{f(x)}$ exists \pause and is equal to $\dfrac{1}{b}$.
  \end{alertblock}
\end{frame}

\begin{frame}
  \frametitle{Limit Properties: Composites}

  \begin{alertblock}{Limits of Composites}\index{limits!of a composite}
    Suppose that $f(x)$ and $g(x)$ are two functions, that $a$ is a real number, and that:
    \begin{enumerate}
      \item $\lim_{x \rightarrow a} f(x)$ exists \pause and is equal to $b$ \pause
      \item $f(x)$ is \textit{not} constantly equal to $b$ around $x = a$, \pause
      \item $\lim_{x \rightarrow b} g(x)$ exists \pause and is equal to $c$ \pause
    \end{enumerate}

    Then we have that $\lim_{x \rightarrow a} g(f(x)) = c$.
  \end{alertblock}
\end{frame}

\begin{frame}
  \frametitle{Combining Continuous Functions}

  We can use the above limit laws to enlarge the class of functions we know to be continuous! \pause

  \begin{alertblock}{Combining Continuous Functions}\index{continuous function!combinations}
    Suppose that $f(x)$ and $g(x)$ are two functions that are continuous at some real number $a$. \pause
    Then we have that: \pause
    \begin{enumerate}
      \item $f(x) + g(x)$ is continuous at $x = a$ \pause
      \item $f(x)g(x)$ is continuous at $x = a$ \pause
      \item If $g(a)$ and $\lim_{x \rightarrow a} g(x)$ are both nonzero then $\dfrac{f(x)}{g(x)}$ is continuous at $x = a$. \pause
      \item If $g(x)$ is instead continuous at $f(a)$ then $g(f(x))$ is continuous at $x = a$.
    \end{enumerate}
  \end{alertblock}
\end{frame}

\begin{frame}
  \frametitle{Moral of the Story}

  The big upshot of this is the following:

  \begin{alertblock}{Strategy for Computing Limits}
    If you are tasked with computing a limit $\lim_{x \rightarrow a} F(x)$ algebraically \pause
    your goal should be to reduce the expression of the functions down \pause
    to end up with an equivalent limit $\lim_{x \rightarrow a} F(x) = \lim_{x \rightarrow a} f(x)$
    where: \pause
    \begin{enumerate}
      \item $f(x)$ is a continuous function \pause
      \item and $a$ is in the domain of $f(x)$.
    \end{enumerate}
  \end{alertblock}
\end{frame}

\begin{frame}
  \frametitle{Moral of the Story: Example}

  For instance:
  \begin{align}
    \lim_{x \rightarrow 1} \dfrac{x^3 - x^2 - x + 1}{x^2 + x - 2} &= \pause \lim_{x \rightarrow 1} \dfrac{x^3 - x^2 - x - 1}{(x - 1)(x + 2)} \pause \\
                                                                  &= \lim_{x \rightarrow 1} \dfrac{(x - 1)(x^2 - 1)}{(x - 1)(x+2)} \pause \\
                                                                  &= \lim_{x \rightarrow 1} \dfrac{(x-1)}{(x-1)} \dfrac{x^2 - 1}{x + 2} \pause \\
                                                                  &= \lim_{x \rightarrow 1} \dfrac{x^2 - 1}{x + 2} \pause \\
                                                                  &= \dfrac{(1)^2 - 1}{1 + 2} \pause = \dfrac{1 - 1}{3} \pause = \dfrac{0}{3} = 0
  \end{align}
\end{frame}

\section{One Sided Limits}

\begin{frame}
  \frametitle{One Sided Limits}

  We can also analyze a similar kind of quantitiy to the limit that will end up being useful when we study limits at infinity on Thursday: \pause

  \begin{block}{One-Sided Limits}\index{limits!one sided}
    Given a function $f(x)$ and a number $a$: \pause
    \begin{itemize}
      \item We say that $\lim_{x \rightarrow a^{-}} f(x) = L$ (verbally, ``the limit of $f(x)$ as $x$ approaches $a$ \textit{from the left}") \pause
        if we can make $f(x) \approx L$ (up to some error $\epsilon > 0$) \pause just by forcing $x \approx a$ (up to some error $\delta > 0$) \pause and $x < a$. \pause
      \item Similarly, we say that $\lim_{x \rightarrow a^{+}} f(x) = L$ (verbally, ``the limit of $f(x)$ as $x$ approaches $a$ \textit{from the right}") \pause
        if we can make $f(x) \approx L$ (up to some error $\epsilon > 0$) \pause just by forcing $x \approx a$ (up to some error $\delta > 0$) \pause and $x > a$.
    \end{itemize}
  \end{block}
\end{frame}

\begin{frame}
  \frametitle{Relationship between Limits and One-Sided Limits}
  
  Questions: \pause
  \begin{enumerate}
    \item Can we have that $\lim_{x \rightarrow a^{-}} f(x) \neq \lim_{x \rightarrow a^{+}} f(x)$? \pause
    \item If $\lim_{x \rightarrow a^{-}} f(x) \neq \lim_{x \rightarrow a^{+}} f(x)$ what can we say about $\lim_{x \rightarrow a} f(x)$? \pause
    \item If $\lim_{x \rightarrow a^{-}} f(x) = \lim_{x \rightarrow a^{+}} f(x)$, what can we say about $\lim_{x \rightarrow a} f(x)$? \pause
    \item If $\lim_{x \rightarrow a^{-}} f(x) = f(a) = \lim_{x \rightarrow a^{+}} f(x)$, what previously discussed definition does the function $f(x)$ and the point $a$ satisfy? \pause
  \end{enumerate}

  \vspace{\stretch{1}}
\end{frame}

\begin{frame}
  \frametitle{Worksheet Time!}

  Let's get into groups and work on the worksheet!
\end{frame}

\begin{frame}
  \frametitle{End of Class!}

  That's all folks!
\end{frame}
\end{document}
